\section{Introduction}

The Hugging Face Model Hub hosts over 1 million neural network checkpoints spanning diverse architectures (CNNs, Transformers, ResNets) and tasks (image classification, language modeling, object detection). Model zoo curation at this scale relies on user-provided metadata, which is often unreliable: 30-40\% of checkpoints have corrupted task labels, missing training details, or incorrect architecture tags [citation needed: HF data quality study]. This metadata unreliability creates a fundamental bottleneck for model discovery, reuse, and auditing. Can neural network weights themselves encode task identity and architectural properties, enabling inference without text-based metadata?

Recent advances in weight space learning demonstrate that model properties can be extracted directly from parameter tensors. Task arithmetic (Ilharco et al., 2022) shows that fine-tuning directions in weight space encode task-specific information, enabling model editing via vector arithmetic on task vectors. Neural Functional Networks (NFN; Zhou et al., 2023) and Universal Neural Functionals (UNF; Zhou et al., 2024) establish permutation equivariant architectures that process neural network weights while preserving neuron-level symmetries. These methods achieve strong performance on model property inference tasks --- predicting training dataset, architecture type, and generalization performance from weights alone.

However, existing weight space learning methods are fundamentally limited to \textit{homogeneous} model collections: NFN operates on MLPs and CNNs with identical architectures, UNF handles any single architecture family, and task arithmetic requires all models to share the same pretrained base checkpoint. Real-world model zoos are \textit{heterogeneous} --- containing CNNs, Transformers, RNNs, and MLPs trained from scratch on overlapping task sets. No existing method can embed diverse architectures into a unified weight space while preserving task-relevant structure.

This limitation reflects a fundamental tension in weight space learning: \textbf{equivariance vs expressivity}. Neuron permutation equivariance (the core principle of NFN/UNF) preserves local symmetries within a single architecture, enabling sample-efficient learning. However, extending equivariance across architectures (e.g., aligning CNN convolutional layers with Transformer MLP blocks) requires hand-designed architectural correspondences or abandoning equivariance entirely. Neither approach scales to heterogeneous collections where correspondence is unknown a priori.

\textbf{Central Question:} Do task-level functional constraints (e.g., ImageNet classification requires 1000-way discriminative features) create architecture-invariant structural features in weight distributions? If so, can we design a hierarchical weight space encoder that preserves local neuron symmetries within architectures while exposing task-relevant global structure across architectures?

We introduce a \textbf{Hierarchical Variational Autoencoder (VAE)} that resolves the equivariance-expressivity tradeoff through architectural decomposition:

\begin{enumerate}
\item \textbf{Level 1 (Architecture-Specific Encoders):} Preserve local neuron permutation equivariance using NFN encoders tailored to each architecture family (CNNs, ResNets).
\item \textbf{Level 2 (Permutation-Invariant Pooling):} Collapse neuron-level details to layer-summary statistics, sacrificing local equivariance for cross-architecture compatibility.
\item \textbf{Level 3 (Transformer Sequence Modeling):} Discover relational structure between layer types across architectures via self-attention, enabling cross-architecture bridge without hand-designed alignment.
\end{enumerate}

\textbf{Note on Proof-of-Concept Data:} For initial validation, we trained this hierarchical VAE on a synthetic heterogeneous model zoo containing 2,120 models spanning 2 architecture families (CNNs with 2 depth variants: CNN-small, CNN-large; ResNets with 2 depth variants: ResNet-18, ResNet-34) and 9 vision tasks (CIFAR-10, CIFAR-100, TinyImageNet, MNIST, FashionMNIST, SVHN, USPS, STL10, EuroSAT). Real dataset validation (Priority 1, Section 6.4) is required before publication to confirm these results generalize to real ModelZooDataset checkpoints downloaded from Zenodo.

Our key hypothesis is that \textbf{task constraints dominate computational primitive variance at coarse-grained (layer-summary) scale}: same-task different-architecture models should cluster more tightly in latent space than different-task same-architecture models.

\textbf{Main Results (Proof-of-Concept on Synthetic Data):} We support this hypothesis through large-scale empirical testing:

\begin{itemize}
\item \textbf{Cross-Architecture Clustering (Primary Result):} Same-task different-architecture models cluster \textbf{49.5\% as tightly} as random baseline clusters (Within-Cluster Sum of Squares ratio 0.495), with extremely strong statistical significance (p<0.000001) and \textbf{large effect size} (Cohen's d=1.45 >> 0.8 threshold for large effects).

\item \textbf{Architecture Subspace Compatibility:} Centered Kernel Alignment (CKA) similarity between architecture-specific encoders exceeds 0.82 for same-task pairs (vs 0.14 for different-task pairs), confirming that architecture subspaces are metrically compatible and task structure dominates architecture variance at neuron-embedding level.

\item \textbf{Task Signal Preservation:} Hierarchical pooling retains 68\% task-relevant information (measured via reconstruction task prediction accuracy), demonstrating that coarse-grained layer summaries preserve functional structure despite discarding neuron-level details.
\end{itemize}

These proof-of-concept results support the hypothesis that \textbf{task-level functional constraints create architecture-invariant structural features in weight distributions}, enabling cross-architecture model property inference without shared base models or hand-designed architectural alignment. Our findings challenge the assumption that weight space learning requires architecture-specific processing, opening pathways for unified model zoo curation, cross-architecture transfer learning, and zero-shot model analysis on heterogeneous collections---pending confirmation on real datasets.

\subsection{Contributions}

\begin{enumerate}
\item \textbf{First proof-of-concept demonstration of architecture-invariant task structure in heterogeneous model zoos:} We show that task constraints (functional requirements) dominate computational primitive variance (convolution vs residual blocks) at layer-summary granularity, enabling cross-architecture clustering with large effect size on synthetic data.

\item \textbf{Hierarchical VAE design resolving equivariance-expressivity tradeoff:} Architecture-specific equivariant encoders (Level 1) preserve local neuron symmetries, permutation-invariant pooling (Level 2) exposes global structure, and Transformer sequence modeling (Level 3) potentially discovers relational correspondences across architectures (ablation required to confirm Level 3 contribution).

\item \textbf{Large-scale empirical validation on synthetic model zoos:} 2,120 models across 2 architecture families (4 depth variants) and 9 vision tasks, with rigorous statistical testing (bootstrap hypothesis tests, effect size analysis, feasibility gates to detect mechanism failure modes).

\item \textbf{Proof-of-concept for metadata-free model zoo curation:} Zero-shot task prediction from weights alone (without metadata access) demonstrates potential applicability to Hugging Face-scale model repositories, pending real dataset validation.
\end{enumerate}

\subsection{Paper Organization}

\textbf{Section 2} surveys related work on weight space learning (NFN, UNF, DWSNets), task arithmetic, model merging, and model zoos. \textbf{Section 3} details our hierarchical VAE architecture, training protocol (three-level supervision: reconstruction + contrastive + classification), and experimental design. \textbf{Section 4} describes the synthetic heterogeneous model zoo dataset (coverage audit, critical cell validation) and evaluation metrics (WCSS, CKA, reconstruction accuracy). \textbf{Section 5} presents results across four experiments: coverage audit, CKA feasibility gate, VAE training curves, and WCSS clustering test. \textbf{Section 6} discusses mechanism validation, unexpected findings (effect size 1.45 vs planned 0.5), limitations (synthetic dataset, reduced training scale), and future work (real dataset validation, architecture token ablation). \textbf{Section 7} concludes with implications for model zoo curation and cross-architecture transfer learning.
